\documentclass[a4paper, serif, 10pt]{moderncv}

%% moderncv
% color and style
\moderncvcolor{black}
\moderncvstyle{banking}

%% page layout/margins
\usepackage[top=2.5cm, bottom=2.5cm, left=1.5cm, right=1.5cm]{geometry}

\nopagenumbers{}

%% Babel config for Dutch language
% ref:
% - https://latex3.github.io/babel/guides/locale-dutch.html
\usepackage[dutch]{babel}

%% fontspec
\usepackage{fontspec}

\IfFontExistsTF{Linux Libertine}{
  \setmainfont[Ligatures={TeX, Common}, Numbers=OldStyle]{Linux Libertine}
}{}
\IfFontExistsTF{Linux Libertine O}{
  \setmainfont[Ligatures={TeX, Common}, Numbers=OldStyle]{Linux Libertine O}
}{}

%% CTeX
% CJK support, used to show CJK characters in the resume
%
% - heading=false: disable ctex heading
% - punct=kaiming: use kaiming punctuations styles for CJK
% - scheme=plain: use plain scheme, do not override \normalsize font size
% - space=auto: space settings for CJK characters
%
% CJK fonts are resolved per locale language, e.g. Noto Serif/Sans CJK JP
% for Japanese resumes, with fallback to the next region if not installed
%
% ref:
% - http://ctan.mirrorcatalogs.com/language/chinese/ctex/ctex.pdf
\usepackage[UTF8, heading=false, punct=kaiming, scheme=plain, space=auto]{ctex}

\IfFontExistsTF{Noto Serif CJK SC}{
  \setCJKmainfont{Noto Serif CJK SC}
}{}
\IfFontExistsTF{Noto Sans CJK SC}{
  \setCJKsansfont{Noto Sans CJK SC}
}{}

%% line spacing
\usepackage{setspace}
\setstretch{1.125}

%% URL styling - use normal text instead of monospace
\urlstyle{same}

%% Patch \httplink and \httpslink to handle full URLs
% The original moderncv macros always prepend http:// or https:// to URLs.
% This causes issues when the URL already has a protocol (e.g., https://example.com)
% resulting in malformed URLs like https://https://example.com.
% This patch checks if the URL already contains "://" and uses it directly if so.
% Note: We use \str_set:Nx (x-expansion) to expand macros like \@homepage first.
\ExplSyntaxOn
\str_new:N \g_yamlresume_url_str

\RenewDocumentCommand{\httpslink}{O{}m}{
  \str_set:Nx \g_yamlresume_url_str {#2}
  \str_if_in:NnTF \g_yamlresume_url_str {://}
    { \url{#2} }
    { \url{https://#2} }
}

\RenewDocumentCommand{\httplink}{O{}m}{
  \str_set:Nx \g_yamlresume_url_str {#2}
  \str_if_in:NnTF \g_yamlresume_url_str {://}
    { \url{#2} }
    { \url{http://#2} }
}
\ExplSyntaxOff

\name{Sanne van Dijk}{}
\title{Senior beveiligingsingenieur met focus op zero trust en incidentrespons}
\phone[mobile]{+31 6 184 27 395}
\email{sanne.vandijk@vandijk-security.nl}
\homepage{https://www.vandijk-security.nl}

\address{Maliebaan 45, Utrecht, Utrecht, Nederland, 3581 CD}{}{}

\extrainfo{{\small \faLinkedin}\ \href{https://www.linkedin.com/in/sannevandijk-cyber}{@sannevandijk-cyber} {} {} {} • {} {} {} 
{\small \faGithub}\ \href{https://github.com/sannevandijk}{@sannevandijk}}

\begin{document}

\maketitle

\section{Basis Informatie}

\cvline{}{\begin{itemize}
\item Ervaren beveiligingsingenieur met meer dan 8 jaar ervaring in het ontwerpen, implementeren en beheren van weerbare IT-omgevingen.
\item Gespecialiseerd in zero-trustarchitectuur, dreigingsdetectie, cloudbeveiliging en incidentrespons voor financiële en zorgorganisaties.
\item Combineert technische diepgang met heldere communicatie richting CISO, auditors en niet-technische stakeholders.
\item Aantoonbaar resultaatgericht in het verlagen van cyberrisico's en het verhogen van de volwassenheid van securityprocessen.
\end{itemize}}
\section{Opleidingen}

\cventry{sep 2014–jul 2017}
        {Master, Cybersecurity, Score: 8.1}
        {Technische Universiteit Delft}
        {\url{https://www.tudelft.nl}}
        {}
        {\begin{itemize}
\item Afstudeeronderzoek naar detectie van supply-chainaanvallen in CI/CD-pipelines.
\item Organiseerde Capture the Flag-evenementen en begeleidde eerstejaars bij securityvakken.
\item Ontwikkelde een prototype voor geautomatiseerde kwetsbaarheidsanalyse in containeromgevingen.
\end{itemize}
\textbf{Cursussen}: Cryptografie, Netwerkbeveiliging, Beveiligingsgovernance, Penetratietesten, Forensische Analyse, Privacy en Gegevensbescherming, Cloudbeveiliging, Risicomanagement}

\cventry{sep 2010–jul 2014}
        {Bachelor, ICT \& Security, Score: 8.4}
        {Hogeschool Utrecht}
        {\url{https://www.hu.nl}}
        {}
        {\begin{itemize}
\item Minor gevolgd in digitale forensica en incidentonderzoek.
\item Stage gelopen bij een regionale zorginstelling en daar een basis-SOC ingericht.
\end{itemize}
\textbf{Cursussen}: Netwerkarchitectuur, Besturingssystemen, Ethisch Hacken, Digitale Forensica, Informatiebeveiliging, Programmeerprincipes}
\section{Werkervaring}

\cventry{apr 2023–Heden}
        {Senior beveiligingsingenieur}
        {Veiligedomeinen B.V.}
        {\url{https://www.veiligedomeinen.nl}}
        {}
        {\begin{itemize}
\item Leidt het zero-trustprogramma voor twaalf bedrijfsonderdelen en adviseert de CISO over risicogebaseerde prioritering.
\item Ontwerpt en implementeert SIEM-detectieregels, SOAR-playbooks en EDR-beleid voor een 24/7 SOC.
\item Voert dreigingsmodellering volgens STRIDE en MITRE ATT\&CK uit voor cloudmigraties naar Azure.
\item Coördineert incidentrespons bij ernstige incidenten, inclusief forensisch onderzoek, crisiscommunicatie en herstelvalidatie.
\item Beheert ISO 27001- en NIS2-compliance en begeleidt interne en externe audits.
\end{itemize}
\textbf{Trefwoorden}: Zero Trust, SIEM, SOAR, Azure-beveiliging, ISO 27001, NIS2, Dreigingsmodellering}

\cventry{jun 2019–mrt 2023}
        {Beveiligingsanalist}
        {Noordzee Fintech}
        {\url{https://www.noordzeefintech.nl}}
        {}
        {\begin{itemize}
\item Bouwde en beheerde detectiecontent in Splunk en Microsoft Sentinel voor betalings- en klantprocessen.
\item Leidde phishing-, malware- en fraudeonderzoeken en vertaalde bevindingen naar concrete maatregelen.
\item Automatiseerde vulnerability management en verkortte de gemiddelde hersteltijd van kritieke kwetsbaarheden met 38\%.
\item Droeg bij aan GDPR- en PSD2-audits en het opstellen van securityrapportages voor het bestuur.
\end{itemize}
\textbf{Trefwoorden}: SOC, Splunk, Microsoft Sentinel, Phishing, Malwareanalyse, GDPR}

\cventry{sep 2017–mei 2019}
        {IT-beveiligingsspecialist}
        {Utrecht Zorgsystemen}
        {\url{https://www.utrechtzorgsystemen.nl}}
        {}
        {\begin{itemize}
\item Beveiligde zorgapplicaties en medische netwerken volgens NEN 7510 en interne beveiligingsrichtlijnen.
\item Implementeerde meerfactorauthenticatie, patchmanagement en netwerksegmentatie voor 3.500 werkplekken.
\item Voerde risicoanalyses uit en ondersteunde teams bij het herstellen van geconstateerde kwetsbaarheden.
\end{itemize}
\textbf{Trefwoorden}: MFA, Patchmanagement, Risicoanalyse, NEN 7510, Firewalls}
\section{Talen}

\cvline{Nederlands}{Moedertaal of tweetalige vaardigheid \hfill \textbf{Trefwoorden}: Moedertaal}
\cvline{Engels}{Volledige professionele vaardigheid \hfill \textbf{Trefwoorden}: TOEFL iBT 112, C1}
\cvline{Duits}{Beperkte werkvaardigheid \hfill \textbf{Trefwoorden}: Goethe B1}
\section{Vaardigheden}

\cvline{Beveiligingsarchitectuur}{Expert \hfill \textbf{Trefwoorden}: Zero Trust, Netwerksegmentatie, Firewalls, VPN, IDS/IPS}
\cvline{Dreigingsdetectie}{Geavanceerd \hfill \textbf{Trefwoorden}: SIEM, Splunk, Microsoft Sentinel, EDR, Dreigingsjacht}
\cvline{Cloudbeveiliging}{Geavanceerd \hfill \textbf{Trefwoorden}: Azure, AWS, IAM, Kubernetes, Terraform}
\cvline{Compliance en Risicobeheer}{Geavanceerd \hfill \textbf{Trefwoorden}: ISO 27001, NIST CSF, GDPR, Risicobeoordeling, Audits}
\cvline{Incidentrespons}{Expert \hfill \textbf{Trefwoorden}: Incidentrespons, Digitale Forensica, Malwareanalyse, SOAR, Crisisbeheersing}
\section{Onderscheidingen}

\cventry{okt 2023}
        {Nederlandse Vereniging voor Cybersecurity}
        {Prijs voor Cybersecurity-Excellentie}
        {}
        {}
        {Uitgereikt voor het succesvol opzetten van een organisatiebreed zero-trustprogramma en het meetbaar verlagen van het aantal kritieke security-incidenten.}
\section{Certificaten}

\cventry{mei 2022}
        {(ISC)²}
        {Certified Information Systems Security Professional (CISSP)}
        {\url{https://www.isc2.org/certifications/cissp}}
        {}
        {}

\cventry{nov 2021}
        {Microsoft}
        {Azure Security Engineer Associate (AZ-500)}
        {\url{https://learn.microsoft.com/en-us/certifications/azure-security-engineer/}}
        {}
        {}

\cventry{mrt 2020}
        {EC-Council}
        {Certified Ethical Hacker (CEH)}
        {\url{https://www.eccouncil.org/programs/certified-ethical-hacker-ceh/}}
        {}
        {}
\section{Publicaties}

\cventry{apr 2024}
        {Zero Trust in de Praktijk: een blauwdruk voor middelgrote organisaties}
        {Cybersecurity Magazine NL}
        {\url{https://www.cybersecuritymagazine.nl/zero-trust-praktijk}}
        {}
        {\begin{itemize}
\item Beschrijft een gefaseerde aanpak voor zero-trustarchitectuur binnen middelgrote organisaties.
\item Gaat in op identiteitsbeheer, netwerksegmentatie, detectie en het meten van securityvolwassenheid.
\end{itemize}}
\section{Projecten}

\cventry{jun 2023–feb 2024}
        {Openbare referentiearchitectuur met beleidsregels, dreigingsmodellen en implementatiechecklisten voor zero trust.}
        {Zero Trust Referentiearchitectuur}
        {\url{https://github.com/sannevandijk/zero-trust-blueprint}}
        {}
        {\begin{itemize}
\item Ontwikkelde een herbruikbare blauwdruk met identity-first principes, microsegmentatie en continue verificatie.
\item Valideerde de architectuur in een Azure-testomgeving met automatische compliancecontroles.
\end{itemize}
\textbf{Trefwoorden}: Zero Trust, Azure, Microsegmentatie, Compliance}

\cventry{feb 2021–dec 2021}
        {Platform dat e-mailheaders, URL-reputatie en gebruikersmeldingen combineert om phishing sneller te detecteren.}
        {Phishingdetectieplatform}
        {\url{https://github.com/sannevandijk/phishing-detection-platform}}
        {}
        {\begin{itemize}
\item Verwerkte dagelijks meer dan 50.000 e-mails en reduceerde valse positieven met 27\%.
\item Integreerde automatische verrijking met dreigingsinformatie en SOAR-playbooks.
\end{itemize}
\textbf{Trefwoorden}: Phishing, Dreigingsinformatie, SOAR, Python}
\section{Interesses}

\cvline{Duursport}{Hardlopen, Wielrennen, Wandelen}
\cvline{Muziek}{Gitaar, Piano}
\cvline{Kennis delen}{Meetups, Mentoring, Capture the Flag}
\section{Vrijwilligerswerk}

\cventry{sep 2021–dec 2024}
        {Vrijwilliger Digitale Weerbaarheid}
        {Veilig Online Utrecht}
        {\url{https://www.veiligonlineutrecht.nl}}
        {}
        {\begin{itemize}
\item Gaf voorlichting aan ouderen en kleine ondernemers over phishing, wachtwoordbeheer en online veiligheid.
\item Organiseerde kwartaalworkshops en stelde laagdrempelig voorlichtingsmateriaal op.
\end{itemize}}
    
\end{document}